\itemhead{第1章 小结与复习 · 第5题（教材p.28）}{XJ-B1-C01-SU-P028-Y1404}
\label{XJ-B1-C01-SU-P028-Y1404}
{\footnotesize 湘教版教材p.28 / PDF第35页 / 第1章 小结与复习 / 第5题}
\field{原题}{数学中的大量命题具有“若 $p$，则 $q$”的形式，试通过对这些命题的梳理，说明充分条件、必要条件以及充要条件的意义。}
\field{参考答案}{$p\Rightarrow q$：$p$ 是 $q$ 的充分条件，$q$ 是 $p$ 的必要条件；双向成立即充要条件。例：$x=1$ 是 $x^2=1$ 的充分而不必要条件。}
{\small \field{关联知识}{L02 充分、必要、充要条件（教材帮p.30、31 / PDF39、40）}
\field{方法与易错点}{分别检查两个方向；用解集包含或具体反例判定，不凭句子位置猜方向。 易错：把“必要”箭头反写；只有一个方向证据便断言“不必要”；忽略给定论域。}
\field{答案依据}{答案册未提供本栏目的独立对应答案；依据原题、教材概念补写参考回应。开放性总结不设唯一数值答案。}}
